\paragraph{Response time, noise, and the sign of tracking.}
For the two-pole response, define the mean delay
$\bar\tau=\gamma^{-1}+\Omega^{-1}$ and
$\mathcal F=\sigma_n^2/(\chi_0/\beta)$. Quantum thermal pairing at
$\beta\hbar=2\pi/H$ implies
\[
\mathcal F\,H\bar\tau\ge4.
\]
The derivation uses the zero-point lower bound on the noise integral
and the logarithmic-mean inequality for the two positive rates.
For the illustrative matter-era occupation at $x=1$, three-percent
tracking requires $H\bar\tau\lesssim0.0253$, hence
$\mathcal F\gtrsim158$. A three-percent force rms then requires
$N_{\rm eff}\gtrsim6.47\times10^4$ under independent-cell averaging.
Spatial correlations, observational windows, and the physical coupling
determine the actual force noise. This inequality constrains the selected
two-pole contact and is not an observational exclusion.

Passive lag also has a definite sign. If an equilibrium target $q(t)$
declines and $\dot n=\Gamma(q-n)$ with $\Gamma>0$ and $n(t_i)\ge q(t_i)$,
then $n(t)\ge q(t)$; positive relaxation cascades retain the inequality.
At fixed baryons, lag behind declining $H$ therefore cannot suppress an
initially equilibrated high-redshift response. Underoccupied assembly
states would require their own preparation history. Relabeling the same
stock between channels leaves its total stress unchanged. In the tested
adiabatic mass-loss model, the surviving mass ratio from $z=2.3$ is
$0.989398$, giving an inferred $a_0$ ratio $m^8=0.9183$, compared with
$H_0/H(2.3)=0.2972$. Neither effect selects the required galactic tracking.

